% Transcription — fonds Alexandre Grothendieck, shelfmark 36, batch 5 (pages 81-84).
% Established from the facsimile archives/batches/36/batch-05.pdf (from a copy
% of https://grothendieck.umontpellier.fr/36.pdf fetched through the project's
% relay), per the skill transcribe-grothendieck. The folder is eighty-four
% pages in five batches; this is the last, four pages, done in parallel with
% the other four.
%
% Pass: Opus 5.5 (claude-opus-5-5), 2026-09-26 — first pass, unchecked against the pages by a human.
%
% Pages skipped: 84, a blank sheet of IHES letterhead with nothing in his hand.
% Page 81 is a tan cover inscribed in ink, top right: a struck-through title
% (« Extensions et biextensions de groupes algébriques ... », the rest lost
% under the strokes, with a boxed « Obstruct[ion?] ») and, below it, the live
% title « Autodualité de la jacobienne, suivant Artin-Mazur. » It takes no
% \page marker; his words become the section heading, with a \note.
%
% Pages summarised: none. Pages transcribed: 82-83, manuscript notes in his
% hand (black ink, fast drafting; the formulas carry, part of the connective
% prose does not).
%
% Typescript: there is none in this batch. The folder's annotated typescript
% (catalogue: « tapuscrit annoté ») lies in the earlier batches; nothing typed
% appears on pages 81-84, so no decision about it is taken here.
% transcripts/36/batch-01.fr.tex did not exist when this batch was finished.
%
% His pagination: none on these leaves; no numbered run. The run is named by
% the cover (p. 81) and by the underlined heading « Cas de dim. relative 1 »
% (p. 82).
%
% What the batch is about. Autoduality of the Jacobian after Artin-Mazur, in
% relative dimension 1: X over a trait S, cohomologically flat in dimension 0;
% P' = Ker(deg : P -> Z_S), Q' = Ker(deg : Q -> Z_S); a 3x3 diagram of exact
% sequences of fppf sheaves with P°, Q°, P', Q', E, E*, Phi, Psi, where
% E, E*, Phi, Psi are represented by étale group schemes with trivial generic
% fibre; on the special fibre Phi = Ker u = Z/dZ (d = pgcd(d_i)), Psi = Coker u
% for u : E(s) -> E*(s), described by lambda = (d_i), mu = (d_i delta_i), alpha,
% beta and the intersection matrix (C_i . C_j)/delta_j. When the local rings of
% X over the strict henselisation are factorial Psi_s is finite; if moreover
% k(s) is perfect or (d^t, p) = 1, Q' is a Néron model of Pic°. Page 83: the
% pairing sum delta_i n_i n'_i on Z^r makes lambda, mu transposed, pairs E(s)
% with E*(s), and u is self-adjoint (symmetry of the intersection form); when
% all delta_i = 1 and d = 1, this yields a canonical autoduality of Psi(s)
% with values in Q/Z, which expresses the obstruction to extending the
% biextension of (P°_eta, P°_eta) by G_m to a biextension of the Néron models
% Q', Q' by G_m.
%
% Notation fixed once: \Phi for his open-bottomed capital phi, \Psi; the
% subscripts « (s) » kept as written (E_{(s)}, E^*_{(s)}) beside the forms
% E(s), E^*(s) he also uses; d^t as written (a small raised t after d, on both
% pages); \underline{n}, \underline{n}' for his underlined vectors.
%
% Readings to check first: « cohom. plat » on p. 82 line 2; the struck
% words above « représentables »; the middle of the N.B. box; « (loc. ...) »
% in the last line of p. 82; « la relation symétrique ... » and the struck
% word before « exprime » on p. 83.

\documentclass[11pt,a4paper]{article}
% Le préambule commun à toutes les transcriptions du fonds Grothendieck.
%
% Il tient en une idée : **l'incertitude se compose, elle ne se résout pas.**
% Une transcription qui lisse un mot douteux a détruit la seule chose pour
% laquelle on la faisait — un lecteur doit pouvoir savoir, à chaque endroit,
% ce qui était sur le feuillet et ce qui a été ajouté. Les six macros qui
% suivent sont donc l'appareil critique, et non de la décoration.
%
% Les mêmes commandes sont reconnues par `scripts/render.mjs`, qui produit la
% vue de lecture du site. Une macro ajoutée ici doit l'être aussi là-bas,
% faute de quoi le rendu échouera bruyamment — ce qui est voulu : un rendu
% silencieusement faux est pire qu'un rendu absent.

\ProvidesPackage{grothendieck}[2026/08/08 Transcriptions du fonds Grothendieck]

\usepackage{amsmath,amssymb,amsthm}
\usepackage{mathrsfs}
\usepackage{stmaryrd}
% Les diagrammes commutatifs sont partout dans ce fonds : c'est la langue
% même de la géométrie algébrique de Grothendieck, et une transcription qui
% les rendrait en prose ne transcrirait rien.
\usepackage{tikz-cd}
% Français par défaut — c'est la langue des feuillets — mais l'anglais est
% chargé aussi : la traduction et le résumé sont des documents anglais, et la
% typographie française (espaces avant les deux-points) y serait une faute.
% Ces documents basculent par \selectlanguage{english} en tête de corps.
\usepackage[english,french]{babel}
\usepackage{fontspec}
\usepackage{geometry}
\geometry{a4paper,margin=25mm,marginparwidth=28mm}
\usepackage{marginnote}
\usepackage{xcolor}
% ulem, not soul. `soul`'s \st and \ul are built on a character-by-character
% scan that XeLaTeX's font machinery breaks: the document fails with an
% undefined control sequence rather than degrading. ulem does the same two jobs
% and is Unicode-engine safe. `normalem` stops it hijacking \emph, which the
% transcriptions use for Grothendieck's own emphasis.
\usepackage[normalem]{ulem}
\usepackage{hyperref}
\hypersetup{colorlinks=true,linkcolor=black,urlcolor=black,pdfborder={0 0 0}}

% --- Métadonnées du lot -----------------------------------------------------
% Reprises telles quelles par le script de rendu, qui les affiche en tête de
% la vue de lecture. Elles fixent ce que le fichier prétend couvrir : sans
% elles, une transcription flotte sans point d'attache dans un fonds de seize
% mille feuillets.

\newcommand{\folder}[1]{\gdef\grothFolder{#1}}
\newcommand{\batch}[1]{\gdef\grothBatch{#1}}
\newcommand{\leaves}[2]{\gdef\grothFirst{#1}\gdef\grothLast{#2}}
\newcommand{\foldertitle}[1]{\gdef\grothFoldertitle{#1}}
\gdef\grothFolder{?}\gdef\grothBatch{?}\gdef\grothFirst{?}\gdef\grothLast{?}\gdef\grothFoldertitle{}

% --- Le résumé --------------------------------------------------------------
%
% Il ouvre la lecture modernisée, sous le titre « Résumé », et il est composé
% en petit. Ce n'est pas un ornement : c'est le seul passage du document qui
% ne soit pas des mathématiques, et un lecteur doit pouvoir le reconnaître
% avant de l'avoir lu — puis le sauter s'il connaît déjà le sujet, ou s'y
% arrêter s'il ne le connaît pas. Le filet qui le suit marque la reprise du
% corps.

\newenvironment{resume}
  {\small\setlength{\parskip}{0.45em}}
  {\par\medskip\hrule\medskip}

% --- Mots-clés ---------------------------------------------------------------
%
% En anglais, en clôture du résumé de la lecture modernisée : le vocabulaire
% moderne sous lequel chercher ce que les feuillets construisent. Le manifeste
% du site (`npm run manifest`) les extrait pour étiqueter la cote entière —
% cette ligne est la source unique des tags, il n'y a pas de fichier à côté.
\newcommand{\keywords}[1]{\par\smallskip\noindent
  {\footnotesize\textsc{Keywords} --- \textit{#1}}\par}

% --- L'appareil critique ----------------------------------------------------

% Le numéro de feuillet, en marge. C'est l'ancre qui permet de régler un
% désaccord sur une formule en regardant *un* feuillet plutôt qu'une chemise
% de six cents. Le site s'en sert aussi pour tourner les pages du fac-similé
% quand on fait défiler la transcription.
\newcommand{\leaf}[1]{%
  \marginnote{\footnotesize\color{gray!70}\textsf{#1}}%
  \label{leaf:#1}\ignorespaces}

% Illisible : on ne devine pas. Un mot inventé qui se lit comme les autres est
% le pire résultat possible.
\newcommand{\ill}{\textcolor{red!60!black}{[\,\ldots\,]}}

% Lecture douteuse : le mot est donné, et signalé comme incertain.
\newcommand{\uncertain}[1]{\uline{#1}}

% Ajout de l'éditeur — une lettre manquante, un mot sous-entendu.
\newcommand{\add}[1]{\textcolor{blue!50!black}{[#1]}}

% Biffé par Grothendieck lui-même. Ce qu'il a rayé fait partie du document :
% une rature dit souvent où la pensée a changé de direction.
\newcommand{\struck}[1]{\sout{#1}}

% Note du transcripteur, sur l'état matériel du feuillet.
\newcommand{\note}[1]{\par\smallskip{\small\color{gray!60!black}\textit{[#1]}}\par\smallskip}

% Note marginale de Grothendieck, distincte du corps du texte.
\newcommand{\marginal}[1]{\par\smallskip\noindent\hspace{1em}%
  {\small\color{gray!60!black}#1}\par\smallskip}

% --- Titre ------------------------------------------------------------------

\AtBeginDocument{%
  \begin{center}
    {\large\bfseries Fonds Alexandre Grothendieck --- cote n\textsuperscript{o}~\grothFolder}\\[2pt]
    {\itshape\grothFoldertitle}\\[4pt]
    {\small feuillets \grothFirst--\grothLast\ (lot \grothBatch)}\\[2pt]
    {\footnotesize\color{gray!60!black}Transcription établie d'après le fac-similé numérisé
     par l'Université de Montpellier.\\
     Les lectures incertaines sont soulignées, les illisibles notées [\,\ldots\,],
     les ajouts entre crochets.}
  \end{center}
  \vspace{1em}}

\endinput


\folder{36}
\batch{5}
\pages{81}{84}
\dating{[vers 1967-1973]}
\watermark{Édition de démonstration}
\foldertitle{SGA 7 (ma part) : notes manuscrites (s.d.), tapuscrit annoté (s.d.)}

\begin{document}

\section*{Autodualité de la jacobienne, suivant Artin-Mazur}

\note{titre inscrit à l'encre en haut à droite du feuillet de garde (page 81),
qui ne porte rien d'autre ; le feuillet ne reçoit pas de numéro de page. Au-dessus,
un premier titre biffé de plusieurs traits obliques : \struck{Extensions et
biextensions de groupes algébriques \ill{} \ill{}}, avec, dans un cadre qui
le rattache, le mot \uncertain{Obstruct}. Pas de pagination de l'auteur sur les
pages 81 à 83.}

\page{82}
\emph{Cas de dim. relative 1}

Supposons \struck{$P \to$} $X$ \uncertain{cohom. plat} en dim.~0, i.e.
$H^0(X_s, \mathcal{O}_{X_s}) \simeq k(s)$.

Posons
\[
P' = \operatorname{Ker}\bigl(P \xrightarrow{\ \deg\ } \mathbb{Z}_S\bigr), \qquad
Q' = \operatorname{Ker}\bigl(Q \xrightarrow{\ \deg\ } \mathbb{Z}_S\bigr).
\]
On a un diagramme commutatif de suites exactes de faisceaux fppf

\begin{tikzcd}[column sep=small, row sep=small]
 & 0 \arrow[d] & 0 \arrow[d] & 0 \arrow[d] & \\
0 \arrow[r] & \Phi \arrow[r] \arrow[d] & P^0 \arrow[r] \arrow[d] & Q^0 \arrow[r] \arrow[d] & 0 \\
0 \arrow[r] & E \arrow[r] \arrow[d] & P' \arrow[r] \arrow[d] & Q' \arrow[r] \arrow[d] & 0 \\
0 \arrow[r] & E/\Phi \arrow[r] \arrow[d] & E^* \arrow[r] \arrow[d] & \Psi \arrow[r] \arrow[d] & 0 \\
 & 0 & 0 & 0 &
\end{tikzcd}

$E, E^*, \Phi, \Psi$ sont \struck{des schémas en g\add{roupes}}
\uncertain{loc. ét.} \uncertain{(quasi-finis)} représentables par des schémas en
groupes \emph{étales}. Leurs fibres génériques sont le groupe unité. Les fibres
spéciales, si $k(s)$ \uncertain{sép. clos}, sont données par
\note{les mots au-dessus de « représentables » sont interlinéaires, ajoutés
au-dessus du passage biffé ; lecture douteuse.}
\[
\Phi_{(s)} = \operatorname{Ker} u \simeq \mathbb{Z}/d\mathbb{Z}, \qquad
\Psi_{(s)} = \operatorname{Coker} u \qquad (d = \operatorname{pgcd}(d_i)),
\]
où
\[
u \colon E_{(s)} \longrightarrow E^*_{(s)}
\]
est l'homom. compris dans le diagramme

\begin{tikzcd}[column sep=huge]
0 \arrow[r] & \mathbb{Z} \arrow[r, "{\lambda=(d_i)}"] & \mathbb{Z}^r \arrow[r, "\alpha"] & E_{(s)} \arrow[r] \arrow[d, dashed, "u"] & 0 \\
 & \mathbb{Z} & \mathbb{Z}^r \arrow[l, "{\mu=(d_i\delta_i)=(d_i^t)}"'] & E^*_{(s)} \arrow[l, "\beta"'] & 0 \arrow[l]
\end{tikzcd}

\note{la flèche $u$ est en pointillé ; devant le $\mathbb{Z}$ de gauche de la
seconde ligne, un « $0 \leftarrow$ » (ou « $0 \to$ ») est biffé.}

et
\[
(\beta u \alpha)\bigl((n_i)_i\bigr) = \Bigl(\sum_{1 \leqslant i \leqslant r}
n_i (C_i \cdot C_j)/\delta_j\Bigr)_{1 \leqslant j \leqslant r} .
\]

N.B. $E(s) \simeq E(S) \to P'(S)$ est donné \uncertain{par son composé avec
$\alpha$, qui est} $(n_i) \mapsto \mathrm{cl}\,\mathcal{O}_X(\sum n_i C_i)$ ;
$P'(S) \to E^*_{(s)} \simeq E^*(S)$ est donné par son composé avec $\beta$, qui
est $L \mapsto \bigl((\deg_{k(s)} L|_{C_i})/\delta_i\bigr)_i$.
\note{ce N.B. est encadré dans la marge droite, sous la ligne « On a, si
$X_{\tilde S}$ est loc. fact. », qui le surmonte.}

On a, si $X_{\tilde S}$ est loc. fact.

\struck{Lorsque $X$ vérifie N, alors $\Psi_{(s)}$ est aussi un gr. fini.}

Si \struck{$k(s)$ est parfait ou $(d^t,p)=1$, \ill{} et} les anneaux locaux
de $X_{\tilde S}$ \add{sont} factoriels, alors $\Psi_s$ est aussi un gr. fini
(comme $\Phi_s$),
\struck{les anneaux locaux de $X_{\tilde S}$ \ill{}}.
Si de plus $k(s)$ est parfait ou $(d^t, p) = 1$, alors
\struck{\ill{} factoriels (pour \ill{}) alors} $Q'$ est un modèle de Néron de
\[
Q'_\eta = P'_\eta = P^0_\eta = \underline{\mathrm{Pic}}^0_{X_\eta/k(\eta)} .
\]
En particulier, si $d^t = 1$, donc $\Phi = 0$, $Q'$ est un modèle de Néron de
$P^0_\eta$, \uncertain{pourvu que} $P'$ \uncertain{soit} (loc. \ill{}) sur $S$
représentable.
\note{les deux dernières lignes sont serrées au pied de la page ; « Néron de
$P^0_\eta$ » est écrit sous la ligne, à gauche. Les biffures des lignes
précédentes montrent l'hypothèse « $k(s)$ parfait ou $(d^t,p)=1$ » déplacée de
la condition de finitude de $\Psi_s$ à celle du modèle de Néron.}

\page{83}
N.B. Accouplons $\mathbb{Z}^r$ et $\mathbb{Z}^r$ par
\[
\bigl((n_i), (n'_i)\bigr) = \sum \delta_i n_i n'_i .
\]
Alors $\lambda$ et $\mu$ sont formellement transposés par cet accouplement,
\struck{\uncertain{en fait}}, donc $\operatorname{Ker}\mu \simeq E^*(s)$ est
accouplé à $\operatorname{Coker}\lambda \simeq E(s)$, par la formule
\[
(\alpha(x), y^*) = (x, \beta y^*) \qquad x \in \mathbb{Z}^r,\ y^* \in E^*(s),
\]
\[
E(s) \times E^*(s) \longrightarrow \mathbb{Z} .
\]
\note{devant $\mathbb{Z}^r$, dans « $x \in \mathbb{Z}^r$ », un premier
$\mathbb{Z}$ est surchargé.}

\struck{Par suite} Je dis que l'on a aussi
\[
(x, u(y)) = (y, u(x)) \qquad \text{pour } x, y \in E(s).
\]
En effet, on peut supposer que $x = \alpha((n_i))$, $y = \alpha((n'_j))$,
\note{sous $(n_i)$ et $(n'_j)$, une accolade renvoie à $\underline{n}$ et
$\underline{n}'$.}
et l'on a donc aussi \uncertain{par} \ill{}
\[
\begin{aligned}
(\alpha(\underline{n}), u(y)) &= (\underline{n}, \beta u \alpha(\underline{n}')) \\
(\alpha(\underline{n}'), u(x)) &= (\underline{n}', \beta u \alpha(\underline{n})) .
\end{aligned}
\]
Or on a
\[
(\underline{n}, \beta u \alpha(\underline{n}')) = \sum n_i n'_j (C_i \cdot C_j)
\]
qui est bien symétrique en $\underline{n}, \underline{n}'$.

Si les $\delta_i$ sont tous $1$ (\uncertain{p.ex.} $k(s)$ parfait), alors on
trouve par l'accouplement précédent
\[
E^*(s) \simeq \operatorname{Hom}(E(s), \mathbb{Z})
\]
et la relation \uncertain{symétrique} \ill{} de plus $d = 1$ i.e.\ $d^t = 1$.
\emph{Dans} ce cas cette dualité, et le fait que $u \colon E(s) \to E^*(s)$
soit \uncertain{autoadjoint}, donne pour le
$\operatorname{Coker} u = \Psi(s)$ une \emph{autodualité} canonique (à valeurs
dans $\mathbb{Q}/\mathbb{Z}$). Cette autodualité \struck{\ill{}} exprime
l'obstruction à étendre la biextension qu'on a en
$(Q'_\eta \times Q'_\eta) = (P^0_\eta \times P^0_\eta)$ par
$\mathbb{G}_m$, en une biextension des $Q', Q'$ par $\mathbb{G}_m$.
\marginal{modèles de Néron}
\note{« modèles de Néron » est entouré sous la ligne, relié par deux traits
aux deux $Q'$. Le feuillet s'arrête ici.}

\end{document}
